\documentclass[10pt,twocolumn]{article}

\usepackage[a4paper,margin=0.72in]{geometry}
\usepackage[T1]{fontenc}
\usepackage[utf8]{inputenc}
\usepackage{lmodern}
\usepackage{microtype}
\usepackage{booktabs}
\usepackage{tabularx}
\usepackage{array}
\usepackage{enumitem}
\usepackage{xcolor}
\usepackage{xurl}
\usepackage{hyperref}
\usepackage{fancyhdr}

\hypersetup{
    bookmarksnumbered=true,
    colorlinks=true,
    linkcolor=black,
    citecolor=black,
    urlcolor=blue
}
\urlstyle{same}
\setlength{\columnsep}{0.28in}
\setlength{\parskip}{0.35em}
\setlength{\parindent}{0pt}
\setlist[itemize]{leftmargin=*,noitemsep,topsep=0.2em}
\setlist[enumerate]{leftmargin=*,noitemsep,topsep=0.2em}
\emergencystretch=2em

\pagestyle{fancy}
\fancyhf{}
\fancyhead[L]{\small Ekza Space Protocol}
\fancyhead[R]{\small Creator-Owned 3D Assets}
\fancyfoot[C]{\small Draft v0.6 -- \thepage}
\renewcommand{\headrulewidth}{0.2pt}
\renewcommand{\footrulewidth}{0pt}

\newcommand{\protocol}[1]{\texttt{#1}}
\newcommand{\bp}{basis points}

\begin{document}

\title{\textbf{Ekza Space Protocol}\\\large A Solana-Native Collaboration Layer for Creator-Owned 3D Assets, Avatars, and Virtual Spaces}
\author{Wotori Movako\\Wotori Studio\\\href{mailto:wotorimovako@gmail.com}{wotorimovako@gmail.com}}
\date{Draft v0.6 \quad May 2026}
\maketitle
\thispagestyle{fancy}

\begin{abstract}
Ekza Space is a protocol stack for collaborative creation, provenance, and monetization of game-ready 3D assets. It is designed for artists, technical artists, writers, animators, world builders, game developers, and studios that need a transparent way to create reusable digital assets while preserving contribution history and ownership logic.

The core thesis is simple: creative production should have a versioned, programmable, and revenue-aware layer similar to what Git introduced for software, but adapted to the realities of 3D production. A final character, prop, environment, rig, animation set, or virtual space is rarely the work of a single person. It may begin as a story fragment, evolve into concept art, become a model, receive UV layout, textures, rigging, animation, metadata, and engine integration. Ekza Space records this chain as a lineage graph and allows finalized releases to carry contributor-share snapshots and revenue vaults.

The current Solana implementation is organized around three protocol elements: \protocol{solana-stellar}, the collaborative asset-universe and release-accounting layer; \protocol{solana-avatars}, the Web3 identity and avatar NFT layer; and \protocol{solana-ekza-space}, the finite virtual-space ownership layer. The most complete implemented flow today is the Stellar-to-Avatars path: a finalized collaborative release can be linked to an avatar drop, and mint fees can be routed back into the Stellar release vault. Together these elements form the foundation for an interoperable creator economy where digital assets can move from collaborative production into live virtual worlds, games, and metaverse interfaces.
\end{abstract}

\section{Disclaimer}
This document is provided for informational and technical discussion purposes only. It does not constitute an offer to sell, a solicitation to buy, or a recommendation regarding securities, tokens, NFTs, or any financial instrument. Nothing in this document should be interpreted as investment, legal, tax, financial, or accounting advice.

The Ekza Space protocol, Wotori Studio interfaces, smart contracts, SDKs, asset pipelines, economic mechanisms, and roadmap are under active development and may change without notice. Any statements about possible future functionality, integrations, revenue models, market demand, adoption, or ecosystem growth are forward-looking and involve risks and uncertainties. Revenue sharing describes a protocol accounting mechanism and does not guarantee income.

Digital assets, NFTs, smart contracts, and blockchain applications involve technical, legal, operational, market, and regulatory risks. Smart contracts may contain bugs or vulnerabilities, off-chain intellectual-property enforcement may require legal processes, and platform availability may depend on third-party infrastructure. Users and partners should perform their own technical, legal, and commercial due diligence before relying on the protocol in production.

\section{Vision}
The internet made creative distribution global, but ownership and production workflows remain highly centralized. The most valuable characters, worlds, animations, and game assets are usually created inside large studios. Contributors may receive salaries or one-time fees, while long-term intellectual-property value, licensing revenue, derivative rights, and platform leverage remain with the studio or publisher.

Ekza Space proposes a different model: a creator-owned production network in which the history of an asset is not erased when the asset becomes commercially useful. Every meaningful contribution can become part of a verifiable production graph. The protocol does not try to replace art direction, studios, or professional pipelines. Instead, it gives creators and teams programmable infrastructure for collaboration, attribution, licensing, and downstream monetization.

The intended result is a new category of creative infrastructure:
\begin{itemize}
    \item a Git-like contribution graph for 3D assets, stories, worlds, rigs, animations, and metadata;
    \item a Solana-native accounting layer for releases, vaults, and contributor revenue claims;
    \item persistent Web3 identities and avatar NFTs that can be used across virtual worlds;
    \item virtual spaces whose settings and access rules are controlled by NFT ownership;
    \item SDKs and interfaces that allow game developers to consume creator-owned assets without rebuilding the ownership layer.
\end{itemize}

\section{Problem Statement}
Modern digital production has three structural problems.

\textbf{First, creative ownership is disconnected from creative labor.} A character or world may depend on concept artists, 3D modelers, texture artists, riggers, animators, technical artists, writers, and tool builders. Once the final product is released, most systems only recognize the publisher or the final asset owner, not the complete chain of contributors.

\textbf{Second, collaboration tools are not revenue-aware.} Git, cloud drives, project managers, and art tools can store files and comments, but they are not designed to represent future licensing rights, revenue shares, or release-level contributor snapshots.

\textbf{Third, NFT metadata alone is too narrow for complex production.} A typical NFT can represent a final asset, but a production-grade 3D avatar or environment contains a deeper structure: source files, optimized GLB/GLTF builds, texture sets, animation clips, rig constraints, usage licenses, compatibility data, and a chain of derivative work. Ekza Space treats the final NFT as only one output of a broader protocol state.

\section{Protocol Thesis}
Ekza Space treats creative work as a sequence of typed, linkable, and releasable artifacts. The protocol records enough state to answer four questions:

\begin{enumerate}
    \item \textbf{What was created?} Asset records point to metadata and files stored through off-chain content-addressed or permanent storage.
    \item \textbf{Who contributed?} Assets are linked to contributors, roles, and upstream parents.
    \item \textbf{What is production-ready?} Draft assets can be reviewed, approved, finalized, and released.
    \item \textbf{How is revenue shared?} A release can store a contributor-share snapshot and route claimable revenue through a vault.
\end{enumerate}

This separates the collaboration layer from the asset-consumption layer. A game may only need a finalized GLB avatar and a license proof. A marketplace may only need an NFT and metadata. A studio may need the full lineage graph. Ekza Space provides a shared protocol foundation while allowing each downstream application to choose how much of that state it uses.

\section{Protocol Stack}
The Ekza Space stack is composed of complementary Solana programs and application layers. The public implementation currently centers on the following repositories: \protocol{solana-stellar}, \protocol{solana-avatars}, and \protocol{solana-ekza-space}. The protocol layer uses Program Derived Addresses (PDAs) for deterministic program-owned state, a standard Solana pattern in which addresses are derived from seeds and the program id, with no private key controlling the PDA \cite{solana_pda}.

\begin{table*}[t]
\centering
\small
\caption{Ekza Space protocol elements}
\begin{tabularx}{\textwidth}{>{\raggedright\arraybackslash}p{0.18\textwidth} >{\raggedright\arraybackslash}p{0.24\textwidth} >{\raggedright\arraybackslash}X >{\raggedright\arraybackslash}p{0.24\textwidth}}
\toprule
\textbf{Layer} & \textbf{Repository / module} & \textbf{Purpose} & \textbf{Primary objects} \\
\midrule
Collaboration and provenance & \protocol{solana-stellar} & Anchor-based protocol for collaborative asset universes. It supports universes, typed assets, asset lineage, approved releases, vaults, and contributor-share snapshots. & Universe, Asset, Lineage, Release, Release Vault, Contributor Share \\
\addlinespace
Identity and avatars & \protocol{solana-avatar} layer; public repo \protocol{solana-avatars} & Web3 identity layer for persistent profiles and avatar NFTs. Profiles act as metaverse passports and can link identity data to selected avatar NFTs. & Profile PDA, Avatar NFT, avatar metadata URI, avatar minter \\
\addlinespace
Virtual spaces & \protocol{solana-ekza-space} & Finite collection of numbered Spaces represented by Metaplex NFTs and PDA-backed settings. Space settings are updateable by NFT owners. & Config PDA, Space PDA, Space NFT, space settings URI \\
\addlinespace
Application and SDK layer & Web app, SDK, metaverse renderer, asset services & User-facing creation, minting, upload, conversion, rendering, and game integration. & UI, SDK methods, IPFS/Arweave manifests, GLB/GLTF assets \\
\bottomrule
\end{tabularx}
\end{table*}

\subsection{Solana Stellar: Collaboration Core}
\protocol{solana-stellar} is the collaboration and accounting heart of the system. It reworks earlier CosmWasm business logic for Solana's account model and is designed around collaborative asset universes. Its README defines the core concepts as Universe, Asset, Lineage, Release, Release Vault, and Contributor Share \cite{solana_stellar}.

A \textbf{Universe} is a collaborative workspace owned by a creator, team, or studio. In the current on-chain program, a Universe stores an owner, open/closed contribution flag, project type, metadata hash, and an immutable collaboration policy. The web interface may present richer creation policies such as closed, approval-required, and open contribution, plus creation prices or review deposits; these are product-layer metadata unless and until they are enforced directly by program constraints or escrow accounts. A \textbf{typed Asset} is an on-chain record for an off-chain artifact such as concept art, a 3D model, a rig, an animation, a script, a texture set, or metadata. \textbf{Lineage} connects assets into a directed acyclic graph, allowing a final asset to reference all upstream work. A \textbf{Release} is a production snapshot of an approved asset. A \textbf{Release Vault} can collect SOL-denominated revenue, while \textbf{Contributor Shares} define allocation in \bp{} independent from the limited creator fields available in ordinary NFT metadata.

This design allows Ekza Space to represent a real production chain:
\begin{center}
\small
Text idea $\rightarrow$ concept art $\rightarrow$ 3D model $\rightarrow$ UV layout $\rightarrow$ textures $\rightarrow$ rig $\rightarrow$ animation $\rightarrow$ release.
\end{center}

The final release can be consumed by a marketplace, avatar minter, game, renderer, or virtual world. In the implemented Stellar-to-Avatars bridge, a finalized Stellar release can initialize avatar-drop state in \protocol{solana-avatars}; the bridge writes Stellar link accounts and marks the Stellar release as linked. The upstream chain remains available for attribution, governance, licensing, auditing, and revenue distribution.

\subsection{Solana Avatar: Identity and Avatar NFTs}
The \protocol{solana-avatar} layer, implemented publicly as \protocol{solana-avatars}, gives users persistent identities for virtual worlds. The repository describes the protocol as a decentralized identity system for the metaverse. It combines a Profile Contract that stores identity data via PDAs and an Avatar Minter Contract that lets creators deploy avatar NFT collections with configurable fees and supply \cite{solana_avatars}.

This layer solves a different problem from \protocol{solana-stellar}. Stellar answers: \emph{How was this asset created, by whom, and how should its release economics work?} Avatar answers: \emph{Who is the user in the virtual world, and which NFT represents them?}

A profile can store or link public identity data such as username, description, and avatar metadata. Avatar NFTs can point to 2D or 3D models through metadata and IPFS links. In the current bridge, finalized Stellar releases can become eligible sources for avatar drops with configurable mint price and supply. When a Stellar-linked avatar is minted, the minting fee can be deposited into the Stellar release vault, while the collaboration accounting remains inside the Stellar layer.

\subsection{Solana Ekza Space: Virtual Space Ownership}
\protocol{solana-ekza-space} is the spatial ownership layer. Its current README describes a minimal Anchor program that manages a finite collection of numbered Spaces. Each Space is represented by a Metaplex NFT plus a PDA with on-chain settings \cite{solana_ekza_space}.

The program contains a Config PDA and Space PDAs. The \protocol{mint\_next\_space} instruction mints the next available Space as a 1/1 NFT, transfers SOL to a treasury, creates the Space PDA, and allows a metadata URI to be provided. The \protocol{update\_space\_settings} instruction updates settings such as name, configuration URI, openness, and editability, with access gated by ownership of the Space NFT.

This gives Ekza Space a place layer: not only characters and assets, but also virtual rooms, galleries, worlds, and collaboration zones whose control is portable and verifiable.

\section{End-to-End Workflow}
An Ekza Space workflow can be understood as a path from identity to creation to release to usage.

\begin{enumerate}
    \item \textbf{Identity.} A creator connects a Solana wallet and creates a profile in the avatar layer. This profile becomes a public creative identity and can later be linked to avatar NFTs or assets.
    \item \textbf{Universe creation.} A team opens a Universe in \protocol{solana-stellar}. The Universe defines the workspace where assets and derivative work will be recorded.
    \item \textbf{Asset contribution.} Contributors add typed assets with metadata URIs. These assets may represent sketches, prompts, scripts, models, rigs, texture files, animations, or engine-ready builds.
    \item \textbf{Lineage linking.} Each new derivative asset can reference its parent assets. This builds a production graph rather than a flat file list.
    \item \textbf{Review and release.} When an asset is ready, its creator submits it and the Universe owner can approve or reject it. Approved assets can be finalized into Releases. The release stores a contributor-share snapshot derived from the Universe collaboration policy.
    \item \textbf{Avatar publishing and integration.} A finalized Stellar release can be published as a \protocol{solana-avatars} avatar drop with mint price and maximum supply. The avatar layer stores link accounts back to the Stellar release, so users and applications can navigate from the minted avatar to its production source.
    \item \textbf{Revenue and claims.} Revenue from direct vault deposits and Stellar-linked avatar minting can be routed into the release vault today. Contributors claim according to the release snapshot. Licensing, rent, sale, Space usage, and downstream application revenue are integration targets that can use the same vault-and-share model.
\end{enumerate}

\section{Data Model}
The protocol deliberately separates on-chain records from off-chain payloads.

\textbf{On-chain state} should contain durable facts that require trust minimization: object identifiers, owner addresses, state transitions, parent links, release status, basis-point shares, vault accounting, and permission checks.

\textbf{Off-chain metadata} should contain large or frequently changing payloads: asset files, GLB/GLTF models, texture maps, preview images, animation manifests, license text, compatibility tags, and editorial descriptions. These payloads can be stored through IPFS, Arweave, or other storage providers, with the on-chain record pointing to the relevant URI.

This separation matters because production-grade 3D content is too large and too diverse to live fully on-chain. The chain should not store every polygon, texture, or animation curve. Instead, it should store the provenance, permissions, and economic state that make the asset trustworthy and interoperable.

\section{NFTs, Metadata, and Why Ekza Extends Them}
NFTs are useful final containers for ownership and transfer, but they are not sufficient by themselves for collaborative production. The Metaplex Token Metadata program is a fundamental Solana program for attaching additional data to fungible and non-fungible tokens. It attaches metadata to mint accounts via PDAs, including creator information, shares, and off-chain JSON URIs \cite{metaplex_token_metadata}.

Ekza Space builds on this idea but expands it in three ways:
\begin{itemize}
    \item \textbf{Lineage beyond the final NFT.} Ekza records upstream creative dependencies, not only the final token metadata.
    \item \textbf{Release-level economics.} Contributor shares can be stored as a release snapshot, independent from metadata creator limits.
    \item \textbf{Application-specific rights.} Asset manifests can describe licensing, game compatibility, derivative permissions, and revenue routes that may not fit into a single NFT metadata object.
\end{itemize}

In this model, the NFT is a portable interface for ownership, identity, or space control. The protocol state behind it gives the asset its production memory.

\section{Contributor Economics}
Ekza Space introduces a programmable approach to collaborative economics. A release can store a contributor-share allocation in \bp{}. When revenue reaches a vault, contributors can claim their allocated portions.

The current \protocol{solana-stellar} implementation supports four collaboration policies:
\begin{itemize}
    \item \textbf{Equal / lineage-equal.} The finalized release computes equal shares across the distinct contributors found in the reachable asset lineage.
    \item \textbf{Contribution weighted.} The finalized release computes shares by counting each contributor's assets in the reachable lineage graph, then normalizing those counts to 10,000 \bp{}.
    \item \textbf{Custom.} The Universe owner provides manual contributor shares before finalization. This mode is intended for explicit off-chain agreements or legacy workflows.
\end{itemize}

Manual share editing is rejected for the automatic lineage policies, and the Universe collaboration policy is immutable after Universe creation. This prevents a team from joining under one economic rule and having that rule silently changed before release.

The model supports several revenue sources:
\begin{itemize}
    \item implemented today: direct SOL deposits into a release vault;
    \item implemented today: Stellar-linked avatar minting fees routed from \protocol{solana-avatars} into a Stellar release vault;
    \item integration target: primary sale of asset NFTs or Space NFTs;
    \item integration target: licensing fees for game or metaverse usage;
    \item integration target: rental or access fees for virtual spaces;
    \item integration target: downstream application revenue routed by integrated partners.
\end{itemize}

The protocol does not require every asset universe to use the same economic policy. A solo creator can receive 100\% through automatic lineage logic if they are the only contributor, or through a custom share. A team can rely on protocol-computed equal or weighted lineage accounting, while a studio can use custom splits for explicit agreements and only release selected assets publicly. The important property is that the release snapshot makes the accounting explicit before revenue is distributed.

\section{Licensing and Intellectual Property}
Blockchain state can record attribution and economic commitments, but it does not automatically solve all intellectual-property questions. Ekza Space should therefore treat licensing as a first-class metadata layer.

A production-ready asset manifest should include:
\begin{itemize}
    \item asset type, version, and release identifier;
    \item author and contributor references;
    \item permitted use cases, such as personal avatar use, game integration, commercial use, remixing, or resale;
    \item restrictions, including prohibited content categories or non-transferable rights;
    \item dependency list and parent asset references;
    \item technical compatibility data, such as polygon budget, rig standard, animation clips, texture resolution, and target runtime.
\end{itemize}

This approach makes Ekza Space more than an NFT minting system. It becomes a rights-aware asset pipeline where final outputs can be safely used by developers who need to know both what an asset is and what they are allowed to do with it.

\section{Interoperability and Developer Experience}
The protocol is designed to be consumed by applications that may not care about every internal detail. A game developer should be able to ask:
\begin{itemize}
    \item Does this wallet own or have a license for this avatar or asset?
    \item Where is the engine-ready GLB/GLTF file?
    \item Which animations and rig conventions are available?
    \item Is this asset approved, finalized, and safe to load?
    \item Which Space does this user own, and what settings does it expose?
\end{itemize}

The SDK layer should abstract these questions into developer-friendly methods. The long-term goal is to let applications integrate Ekza assets without becoming blockchain specialists. A WebGL renderer, Bevy game, Three.js application, or custom engine should be able to fetch profiles, avatars, spaces, and releases through stable APIs while still relying on Solana for ownership and accounting.

\section{Security Model}
The protocol security model depends on a clear division of responsibility.

\textbf{Program authority.} PDAs provide deterministic, program-controlled accounts. Only the owning program can sign for a PDA through Solana runtime mechanisms, which is why PDAs are suitable for vaults, profiles, spaces, and other program-owned state \cite{solana_pda_accounts}.

\textbf{Ownership gating.} Some actions, such as updating Space settings, should require proof that the signer owns the relevant NFT token account. This makes ownership portable while keeping settings under protocol control.

\textbf{State transitions.} Asset states are explicit: draft, submitted, approved, rejected, finalized, minted, and archived. Release states are explicit: draft, finalized, linked, and archived. Once a release is finalized, its contributor-share snapshot should not be silently rewritten.

\textbf{Testing and auditability.} The public repositories include Anchor tests or test instructions, including LiteSVM usage for \protocol{solana-ekza-space} and \protocol{anchor test} for \protocol{solana-stellar}. Before production launch, the protocol should undergo expanded tests, adversarial review, and independent smart-contract audits.

\section{Governance and Moderation}
Open creation does not mean absence of rules. Ekza Space needs moderation and governance mechanisms at several levels:
\begin{itemize}
    \item workspace-level rules for who may contribute to a Universe;
    \item approval workflows before an asset becomes a release;
    \item content policy enforcement for public discovery, marketplace listings, and featured spaces;
    \item dispute processes for contributor shares, copied work, and licensing violations;
    \item administrative recovery or freeze policies, if any, disclosed clearly to users.
\end{itemize}

The current on-chain review authority is the Universe owner: creators can submit assets, and the owner can approve or reject them. Richer admin sets, paid creation gates, refundable review deposits, moderation queues, and dispute workflows should be treated as product or roadmap features unless encoded directly in the program. The protocol should prefer transparent rules and narrow administrative powers. Where off-chain decisions are unavoidable, they should be separated from immutable on-chain facts and documented in the user interface.

\section{Current Development Status}
Ekza Space is in an early protocol and product phase. Public repositories show three working directions:
\begin{itemize}
    \item \protocol{solana-stellar}: a Solana protocol core for universes, assets, lineage, releases, contributor shares, vault deposits, and revenue claims \cite{solana_stellar};
    \item \protocol{solana-avatars}: an avatar identity and NFT minting system with profile PDAs, avatar drops, metadata links, a browser interface, and Stellar-linked avatar-drop support \cite{solana_avatars};
    \item \protocol{solana-ekza-space}: a minimal Anchor program for finite numbered Spaces, 1/1 NFTs, treasury payments, and ownership-gated settings \cite{solana_ekza_space}.
\end{itemize}

Wotori Studio positions Ekza Space as a collaboration and monetization ecosystem for 3D assets, described as a GitHub-like system for 3D artists where contributors co-build reusable assets while keeping ownership of their work \cite{wotori_site}.

\section{Implementation Boundaries}
Several product concepts are intentionally broader than the current smart-contract surface.

\textbf{Creation policy and deposits.} The program currently enforces a binary Universe contribution gate: open Universes allow public asset creation, while closed Universes restrict creation to the owner. The web interface may expose closed, approval-required, and open policies, entity creation price, and review deposit fields through metadata. For these to become trust-minimized economic rules, they should be represented as explicit on-chain fields, escrow movements, and refund/slash instructions.

\textbf{Avatar publishing.} The implemented bridge publishes a finalized Stellar release as avatar-drop state in \protocol{solana-avatars}; it does not automatically attach the avatar to a profile or place it inside a Space. Those are downstream integration steps.

\textbf{Licensing and manifests.} License kind is represented on assets, and richer license text, compatibility data, file manifests, and derivative permissions should live in signed or content-addressed metadata. The protocol can reference and enforce parts of that metadata over time, but the current implementation should not be described as a full legal-rights enforcement engine.

\section{Roadmap}
The roadmap can be organized into five phases.

\textbf{Phase 1: Protocol hardening.} Complete unit and integration tests, improve account constraints, finalize state machines, add error coverage, and document all instructions.

\textbf{Phase 2: Asset manifest standard.} Define a stable JSON schema for 3D assets, avatars, spaces, licenses, compatibility data, and contributor metadata.

\textbf{Phase 3: Creator interface.} Build a production-ready web interface for creating universes, uploading assets, linking lineage, approving releases, minting avatars, and managing spaces.

\textbf{Phase 4: SDK and engine integrations.} Provide SDK methods for profile lookup, ownership verification, asset fetching, release validation, and space loading. Prioritize WebGL/Three.js and game-engine friendly GLB/GLTF flows.

\textbf{Phase 5: Marketplace and partner ecosystem.} Enable licensing, rental, revenue routing, creator discovery, public collections, and downstream integrations with games and virtual worlds.

\section{Risks and Limitations}
Ekza Space should be evaluated with realistic expectations.

\textbf{Technical risk.} Smart contracts can contain vulnerabilities, and off-chain storage can fail or become unavailable if not designed carefully.

\textbf{Legal risk.} On-chain attribution does not automatically resolve copyright, trademark, moral rights, work-for-hire arrangements, or jurisdiction-specific disputes.

\textbf{Adoption risk.} A creator-owned economy depends on artists, developers, players, and platforms choosing to use the same asset and identity standards.

\textbf{Economic risk.} NFT sales, licensing demand, and marketplace liquidity are uncertain. Contributor-share mechanisms only distribute revenue that actually exists.

\textbf{Content risk.} Open asset creation requires moderation, safety policies, and dispute handling to prevent plagiarism, abuse, illegal content, or brand misuse.

\section{Conclusion}
Ekza Space is an attempt to give creative collaboration a protocol layer. It brings together asset lineage, contributor attribution, release snapshots, vault-based revenue claims, Web3 avatar identity, and virtual space ownership into a single Solana-native ecosystem.

The core idea is not merely to mint NFTs. The deeper goal is to preserve the memory of creation: who imagined the character, who shaped the model, who painted the textures, who rigged the body, who animated the movement, who made it usable in a world, and how each participant can remain connected to the asset after it becomes valuable.

If successful, Ekza Space can become infrastructure for a new visual economy: one where independent creators can collaborate at studio scale without surrendering the entire future of their work, and where games and virtual worlds can load creator-owned assets with clear provenance, identity, and licensing.

\section{Ecosystem Websites}
\begin{itemize}
    \item Wotori Studio: \url{https://wotori.io}
    \item Ekza Space: \url{https://ekza.io}
    \item Avatar interface: \url{https://avatar.ekza.io}
\end{itemize}

\section{Contacts}
Email: \href{mailto:wotorimovako@gmail.com}{wotorimovako@gmail.com}

\small
\raggedright
\begin{thebibliography}{9}
\bibitem{solana_stellar} Ekza Space, \emph{solana-stellar README}. \href{https://github.com/ekza-space/solana-stellar}{github.com/ekza-space/solana-stellar}
\bibitem{solana_avatars} Ekza Space, \emph{solana-avatars README}. \href{https://github.com/ekza-space/solana-avatars}{github.com/ekza-space/solana-avatars}
\bibitem{solana_ekza_space} Ekza Space, \emph{solana-ekza-space README}. \href{https://github.com/ekza-space/solana-ekza-space}{github.com/ekza-space/solana-ekza-space}
\bibitem{wotori_site} Wotori Studio, \emph{Wotori Studio website}. \href{https://www.wotori.io/}{wotori.io}
\bibitem{solana_pda} Solana Foundation, \emph{Program Derived Addresses}. \href{https://solana.com/docs/core/pda}{solana.com/docs/core/pda}
\bibitem{solana_pda_accounts} Solana Foundation, \emph{PDA Accounts}. \href{https://solana.com/docs/core/pda/pda-accounts}{solana.com/docs/core/pda/pda-accounts}
\bibitem{metaplex_token_metadata} Metaplex, \emph{Token Metadata Overview}. \href{https://www.metaplex.com/docs/smart-contracts/token-metadata}{metaplex.com/docs/smart-contracts/token-metadata}
\end{thebibliography}

\end{document}
